\section*{Hoofdstuk \ref{ch:b2:asexual-reproduction} --- Ongeslachtelijke voortplanting en klonen bij planten}

\begin{solution}{exo:b2:asexual-reproduction:1}
\omterm{def:b2:asexual-reproduction:clone}{Kloon}: de genetisch identieke nakomelingen van één ouder door mitose.
\omterm{def:b2:asexual-reproduction:clone}{Genet}: het genetische individu, alles wat uit één zygote is gegroeid.
\omterm{def:b2:asexual-reproduction:clone}{Ramet}: een fysiologisch zelfstandige eenheid van een \omterm{def:b2:asexual-reproduction:clone}{genet}. Organen:
uitlopers (aardbei), wortelstokken (kweekgras), knollen (aardappel), bollen
(ui), wortelopslag (ratelpopulier), plantjes op bladeren
(\emph{Kalanchoe}).
\end{solution}

\begin{solution}{exo:b2:asexual-reproduction:2}
Het cambium van ent en onderstam moet op elkaar aansluiten en contact
maken; callus overbrugt de snede en differentieert er xyleem en floëem
doorheen. De onderstam levert de wortels (water, mineralen, groeikracht,
tolerantie voor de bodem, beheersing van de grootte), de ent de
vruchtdragende scheut (het ras). Elk deel behoudt zijn eigen genoom en de
cellen versmelten nooit: twee \omterm{def:b2:meiosis-heredity:vocabulary}{genotypen} naast elkaar, een chimaera; een
hybride zou in elke cel een mengsel van beide genomen dragen.
\end{solution}

\begin{solution}{exo:b2:asexual-reproduction:3}
\omterm{def:b2:asexual-reproduction:clone}{Vegetatieve vermeerdering}: een lichaamsdeel schiet wortel en wordt een
nieuwe plant. \omterm{def:b2:asexual-reproduction:apomixis}{Apomixie}: een \omterm{def:b2:angiosperm-reproduction:seed}{zaad} ontstaat zonder bevruchting, met een
embryo dat een \omterm{def:b2:asexual-reproduction:clone}{kloon} van de moeder is. \omterm{def:b2:asexual-reproduction:apomixis}{Parthenogenese}: een dier ontwikkelt
zich uit een onbevruchte eicel. Alleen \omterm{def:b2:asexual-reproduction:apomixis}{apomixie} behoudt het \omterm{def:b2:angiosperm-reproduction:seed}{zaad} --- met
zijn \omterm{def:b2:angiosperm-reproduction:seed}{zaadhuid}, \omterm{def:b2:angiosperm-reproduction:seed}{kiemrust} en verspreiding.
\end{solution}

\begin{solution}{exo:b2:asexual-reproduction:4}
Steriliseer het explantaat; plaats het op agar met suiker, mineralen,
vitaminen en auxine en cytokinine in evenwicht (callus); verhoog de
verhouding van cytokinine tot auxine om scheuten te induceren; zet de
scheuten om de paar weken over; breng ze over naar een auxinerijk medium
om te wortelen; laat de plantjes wennen bij hoge luchtvochtigheid en
daarna in grond.
\end{solution}

\begin{solution}{exo:b2:asexual-reproduction:5}
Elke plant wordt vervangen door $1 + 24 = 25$: na 4 jaar $25^{4} =
\num{390625}$ planten, die \qty{39000}{m^2} nodig hebben, bijna vier
hectare. Ruimte, licht en voedingsstoffen, het korte bereik van een
uitloper en het afsterven van planten bij verdringing beëindigen de
exponentiële fase binnen enkele jaren; de \omterm{def:b2:asexual-reproduction:clone}{kloon} rukt daarna als een front
op.
\end{solution}

\begin{solution}{exo:b2:asexual-reproduction:6}
$\num{14000}/130 \approx 110$ generaties stammen. Straal van een schijf
van \qty{4.3e5}{m^2}: $\sqrt{4.3\times 10^{5}/\pi} = \qty{370}{m}$;
opmars $370/\num{14000} = \qty{0.026}{m}$ per jaar, minder dan drie
centimeter.
\end{solution}

\begin{solution}{exo:b2:asexual-reproduction:7}
$p_t = 2^{t}p_0/(1 + (2^{t} - 1)p_0)$: $t = 5$: $0.0032$; $t = 10$:
$0.1024/1.1023 = 0.093$; $t = 15$: $3.28/4.28 = 0.77$. Hij verspreidt
zich nooit als zijn fitness lager is dan de helft van die van de
geslachtelijke lijn: een relatieve fitness $w$ geeft $q' = 2wq$, dat
afneemt voor $w <
\tfrac12$.
\end{solution}

\begin{solution}{exo:b2:asexual-reproduction:8}
$10 : 1$: wortels; $1 : 1$: ongedifferentieerd callus; $1 : 10$: scheuten.
Een stek heeft al een scheut en heeft wortels nodig, die auxine aan de
afgesneden basis induceert.
\end{solution}

\begin{solution}{exo:b2:asexual-reproduction:9}
1, 10, 100, \num{1000}, \num{10000} kg: na het derde seizoen pas
\qty{1}{t}, na het vierde \qty{10}{t}: vier seizoenen. Echt \omterm{def:b2:angiosperm-reproduction:seed}{zaad} geeft
duizenden zaden per plant, maar de aardappel is een heterozygote
tetraploïde, zodat de zaailingen uitsplitsen en geen enkele het ras is.
\end{solution}

\begin{solution}{exo:b2:asexual-reproduction:10}
$e^{-5} = 0.0067$: 670 mutatievrije individuen in de grote populatie,
ongeveer 7 in de kleine. Zeven individuen kunnen binnen enkele generaties
allemaal door toeval zonder nakomelingen blijven (de kans dat geen van hen
zich voortplant, is per generatie van de orde $e^{-7}$), waarna de beste
klasse voorgoed verdwenen is en de op één na beste de beste wordt: een
klik van de ratel, herhaald tot de last ondraaglijk is. Met 670 is het
verlies praktisch onmogelijk --- tot een flessenhals $N$ doet krimpen of
de \omterm{thm:b2:genome-diversification:rate}{mutatiesnelheid} stijgt; de ratel draait maar in één richting.
\end{solution}

\begin{solution}{exo:b2:asexual-reproduction:11}
In de \omterm{def:b2:asexual-reproduction:clone}{kloon} is elke gastheer vatbaar: de stam verspreidt zich ongehinderd
in het tempo waarin de sporen reizen. In de geslachtelijke populatie is het
\omterm{def:b2:meiosis-heredity:vocabulary}{genotype} dat hij doorbreekt een van $2^{10} = 1024$ combinaties, aanwezig
in ongeveer een duizendste van de planten, verspreid tussen resistente
buren, zodat de epidemie uithongert. De Lumper was één \omterm{def:b2:asexual-reproduction:clone}{kloon}, geteeld over
het grootste deel van een land; de aardappelziekte vond elke plant
onbeschermd, en er was geen variatie om mee te veredelen tot er andere
rassen werden ingevoerd.
\end{solution}

\begin{solution}{exo:b2:asexual-reproduction:12}
De \omterm{prop:b2:asexual-reproduction:whysex}{ratel van Muller}, het combineren van gunstige \omterm{def:b2:genome-diversification:mutation}{mutaties} en de Rode
Koningin geven seks elk een voordeel dat in een grote, geparasiteerde,
veranderende wereld een factor twee kan overtreffen. Bdelloïde
raderdiertjes zijn het lastige geval: veertig miljoen jaar zonder
mannetjes en geen ineenstorting door de ratel. Ze lijken te ontsnappen
door een extreme bestandheid tegen uitdroging (die hun parasieten doodt),
door vreemd DNA op te nemen en door een ongewoon efficiënt \omterm{prop:b2:genome-diversification:repair}{DNA-herstel} ---
uitzonderingen die laten zien wat de regel normaal vereist.
\end{solution}

\begin{solution}{pb:b2:asexual-reproduction:1}
\textbf{1.} $4\times 10^{4}$ planten.
\textbf{2.} $1 + 6\times 2 = 13$.
\textbf{3.} \num{1300}; \num{16900}; \num{219700}.
\textbf{4.} $13^{t} = 400$: $t = \ln 400/\ln 13 = 5.99/2.56 = 2.3$
jaar --- vol in de loop van het derde jaar.
\textbf{5.} Straal van een schijf van een hectare $\sqrt{10^{4}/\pi} = \qty{56}{m}$;
bij \qty{0.5}{m} per jaar \num{113} jaar.
\textbf{6.} Honderd verspreide stichters vullen het veld exponentieel
vanuit honderd fronten in twee jaar; één stichter zou zijn omgeving in
twee jaar vullen en daarna een eeuw lang voortkruipen.
\textbf{7.} $5^{n} \ge 10^{6}$: $n \ge 8.6$, negen overzettingen
($1.95\times 10^{6}$ scheuten), 36~weken.
\textbf{8.} $0.8\times 1.95\times 10^{6} = 1.6\times 10^{6}$ planten.
\textbf{9.} $0.98$; $0.98^{10} = 0.82$.
\textbf{10.} $10\times 0.98 = 9.8$ schone lijnen verwacht; uit tien
stekken één.
\textbf{11.} Het virus verplaatst zich via plasmodesmata en het floëem, en
de apicale koepel heeft geen vaatverbinding en deelt sneller dan het virus
zich verspreidt, zodat het buitenste tiende van een millimeter meestal
niet besmet is.
\textbf{12.} Laboratorium: meteen uniforme, virusvrije planten, maar duur,
en elke \omterm{def:b2:genome-diversification:mutation}{mutatie} of besmetting in de kweek wordt een miljoen keer
vermenigvuldigd. Uitlopers: gratis en ter plaatse, maar drie seizoenen
verloren en elk virus van de moederplanten gaat mee.
\textbf{13.} Apomicten zijn met $pN$ en geven elk $F$ zaden; geslachtelijke
planten zijn met $(1 - p)N$ en geven elk $F/2$ zaden (de helft van hun
inspanning is \omterm{def:b2:plant-life-cycles:heterospory}{stuifmeel}, dat geen zaden maakt); het aandeel van de
apomicten in de zaden is $Fp/(Fp + F(1 - p)/2)
= 2p/(1 + p)$.
\textbf{14.} $p_5 = 0.032/1.031 = 0.031$; $p_{10} = 1.024/2.023 =
0.51$: de frequentie passeert de helft in generatie 10.
\textbf{15.} $p' = 2p(1 - cp)/\bigl(2p(1 - cp) + (1 - p)\bigr)$.
\textbf{16.} $p' = p$ wanneer $2(1 - cp^*) = 1$: $p^* = 1/2c = 0.625$.
Onder $p^*$ is $2(1 - cp) > 1$ en stijgt $p$; erboven daalt ze: stabiel.
\textbf{17.} $p^* \ge 1$ wanneer $c \le \tfrac12$: de apomict neemt het
over, tenzij de parasiet hem, wanneer hij algemeen is, meer dan de helft
van zijn productie kost.
\textbf{18.} $c = 0.6$: $p^* = 0.83$. Om seks in de populatie te houden,
moet de parasiet bij een hoge frequentie van de apomict de fitness van de
\omterm{def:b2:asexual-reproduction:clone}{kloon} met meer dan het tweevoudige voordeel verlagen --- de Rode Koningin
moet hard lopen.
\textbf{19.} Drie chromosoomsets kunnen bij de \omterm{def:b2:meiosis-heredity:meiosis}{meiose} niet paren, zodat zijn
gameten ongebalanceerd zijn; hij is buitengesloten van recombinatie, en de
\omterm{prop:b2:asexual-reproduction:whysex}{ratel van Muller} en de parasieten werken ongehinderd op hem in --- en
daarom zijn apomictische lijnen van paardenbloemen jong en ontstaan ze
steeds opnieuw uit geslachtelijke voorouders.
\textbf{20.} $10^{-9}\times 4.8\times 10^{8} = 0.48$ per deling.
\textbf{21.} $2\times \num{14000}\times 20 = 5.6\times 10^{5}$
delingen; $0.48\times 5.6\times 10^{5} = 2.7\times 10^{5}$
\omterm{def:b2:genome-diversification:mutation}{mutaties}.
\textbf{22.} $2.7\times 10^{5}/4.8\times 10^{8} = 5.6\times 10^{-4}$ van
het genoom; ongeveer \num{5400} coderende verschillen.
\textbf{23.} Competitie tussen cellijnen in het meristeem (een mutante lijn
die trager deelt, wordt verdrongen) en tussen \omterm{def:b2:asexual-reproduction:clone}{rameten} (een stam met een
slechte \omterm{def:b2:genome-diversification:mutation}{mutatie} geeft minder wortelopslag). Zwakker omdat de meeste
somatische \omterm{def:b2:genome-diversification:mutation}{mutaties} \omterm{def:b2:meiosis-heredity:vocabulary}{heterozygoot} en gemaskeerd zijn, en niets ze blootlegt:
geen \omterm{def:b2:meiosis-heredity:meiosis}{meiose} maakt ze \omterm{def:b2:meiosis-heredity:vocabulary}{homozygoot} en geen ééncellig zygotestadium test het
hele genoom in één cel.
\textbf{24.} Niet uniform: zijn stammen verschillen op honderdduizenden
plaatsen. Maar de verschillen zijn grotendeels neutraal en gemaskeerd,
zodat de \omterm{def:b2:asexual-reproduction:clone}{kloon} in feite één \omterm{def:b2:meiosis-heredity:vocabulary}{genotype} is --- een mozaïek van somatische
varianten van één \omterm{def:b2:asexual-reproduction:clone}{genet}.
\textbf{25.} Veld vol na 2.3 jaar; één meristeem geeft in negen maanden
$1.6\times 10^{6}$ planten; met $c = 0.8$ komt de apomict tot rust bij
$p^* = 0.625$, terwijl de geslachtelijke planten \qty{37.5}{\%} behouden.
\end{solution}
